\documentclass{beamer}
\usepackage[utf8]{inputenc}

\usetheme{Madrid}
\usecolortheme{default}
\usepackage{amsmath,amssymb,amsfonts,amsthm}
\usepackage{txfonts}
\usepackage{tkz-euclide}
\usepackage{listings}
\usepackage{adjustbox}
\usepackage{array}
\usepackage{tabularx}
\usepackage{gvv}
\usepackage{lmodern}
\usepackage{circuitikz}
\usepackage{tikz}
\usepackage{graphicx}
\usepackage{multicol}

\setbeamertemplate{page number in head/foot}[totalframenumber]

\usepackage{tcolorbox}
\tcbuselibrary{minted,breakable,xparse,skins}



\definecolor{bg}{gray}{0.95}
\DeclareTCBListing{mintedbox}{O{}m!O{}}{%
  breakable=true,
  listing engine=minted,
  listing only,
  minted language=#2,
  minted style=default,
  minted options={%
    linenos,
    gobble=0,
    breaklines=true,
    breakafter=,,
    fontsize=\small,
    numbersep=8pt,
    #1},
  boxsep=0pt,
  left skip=0pt,
  right skip=0pt,
  left=25pt,
  right=0pt,
  top=3pt,
  bottom=3pt,
  arc=5pt,
  leftrule=0pt,
  rightrule=0pt,
  bottomrule=2pt,
  toprule=2pt,
  colback=bg,
  colframe=orange!70,
  enhanced,
  overlay={%
    \begin{tcbclipinterior}
    \fill[orange!20!white] (frame.south west) rectangle ([xshift=20pt]frame.north west);
    \end{tcbclipinterior}},
  #3,
}
\lstset{
    language=C,
    basicstyle=\ttfamily\small,
    keywordstyle=\color{blue},
    stringstyle=\color{orange},
    commentstyle=\color{green!60!black},
    numbers=left,
    numberstyle=\tiny\color{gray},
    breaklines=true,
    showstringspaces=false,
}
%------------------------------------------------------------
%This block of code defines the information to appear in the
%Title page
\title %optional
{8.4.32}
\date{October 12,2025}


\author 
{Jnanesh Sathisha karmar - EE25BTECH11029}



\begin{document}



\frame{\titlepage}
\begin{frame}{Question:}


The line passing through the extremity \textbf{A} of the major axis and extremity \textbf{B} of the minor axis of the ellipse $x^2 + 9y^2 = 9$ meets its auxiliary circle at the point \textbf{M}. Then the area of the triangle with the vertices at \textbf{A}, \textbf{M}, and the origin \textbf{O} is (2009) \\
\begin{enumerate}
\begin{multicols}{4}
    \item $\frac{31}{10}$
    \item $\frac{29}{10}$ 
    \item $\frac{21}{10}$
    \item  $\frac{27}{10}$
\end{multicols}
\end{enumerate}
\end{frame}

\begin{frame}{Theoretical Solution}

\noindent\textbf{Solution}\\
The solution is derived using a linear algebra approach by representing the geometric elements as vectors and matrices.

The given equation of the ellipse is $x^2 + 9y^2 = 9$. We can express this in the standard matrix form for a conic, $\vec{x}^\top\vec{V}\vec{x} + f = 0$.
Let the position vector $\vec{x}$ be defined as:
\begin{align}
    \vec{x} = \myvec{x \\ y}
\end{align}
    

The equation can then be written as:
\begin{align}
     \vec{x}^\top \myvec{1 & 0 \\ 0 & 9} \vec{x} - 9 = 0
\end{align}
   
Here, the matrix $\vec{V} = \myvec{1 & 0 \\ 0 & 9}$ and the scalar $f = -9$.To find the semi-axes, we normalize the equation to the form $\vec{x}^\top\vec{A}\vec{x} = 1$:

\end{frame}
\begin{frame}{Theoretical Solution}
\begin{align}
    \frac{x^2}{9} + \frac{y^2}{1} = 1 \implies \vec{x}^\top \myvec{1/9 & 0 \\ 0 & 1} \vec{x} = 1
\end{align}
The semi-axis lengths are the reciprocal of the square roots of the eigenvalues of the matrix $\vec{A} = \myvec{1/9 & 0 \\ 0 & 1}$.
\begin{itemize}
    \item The eigenvalues are $\lambda_1 = 1/9$ and $\lambda_2 = 1$.
    \item The major semi-axis is $a = 1/\sqrt{\lambda_1} = 3$. The corresponding eigenvector is $\vec{v}_a = \myvec{1 \\ 0}$.
    \item The minor semi-axis is $b = 1/\sqrt{\lambda_2} = 1$. The corresponding eigenvector is $\vec{v}_b = \myvec{0 \\ 1}$.
\end{itemize}
The vertices $\vec{A}$ and $\vec{B}$ are at the extremities of these axes.

\end{frame}
\begin{frame}{Theoretical Solution}
\begin{itemize}
    \item Position vector of $\vec{A}:$ $\vec{a} = a \cdot \vec{v}_a = 3 \myvec{1 \\ 0} = \myvec{3 \\ 0}$. So, A = (3, 0).
    \item Position vector of \textbf{B}: $\vec{b} = b \cdot \vec{v}_b = 1 \myvec{0 \\ 1} = \myvec{0 \\ 1}$. So, B = (0, 1).
\end{itemize}
The vector equation of the line passing through $\vec{A}$ and $\vec{B} $is:
\begin{align}
    \vec{x}\brak{t} = \vec{a} + t\brak{\vec{b} - \vec{a}} = \myvec{3 \\ 0} + t\brak{\myvec{0 \\ 1} - \myvec{3 \\ 0}} = \myvec{3 - 3t \\ t}
\end{align}
The auxiliary circle has the equation $x^2 + y^2 = a^2$, which is $x^2 + y^2 = 9$. We find the intersection point \textbf{M} by substituting the parametric equations of the line into the circle's equation:
\begin{align}
 \brak{3 - 3t}^2 + t^2 = 9\\
    9 - 18t + 9t^2 + t^2 = 9\\
    10t^2 - 18t = 0 \implies 2t\brak{5t - 9} = 0
\end{align}

\end{frame}
\begin{frame}{Theoretical Solution}
The solutions are $t=0$ (which corresponds to point $\vec{A}$) and $t = 9/5$. The position vector of $\vec{M}$ is found using $t = 9/5$:
\begin{align}
     \vec{m} = \myvec{3 - 3\brak{9/5} \\ 9/5} = \myvec{15/5 - 27/5 \\ 9/5} = \myvec{-12/5 \\ 9/5}
\end{align}
   
The area of a triangle with vertices at the origin \textbf{O}, \textbf{A}, and \textbf{M} is given by half the absolute value of the determinant of the matrix formed by the position vectors $\vec{a}$ and $\vec{m}$.
\[
    \text{Area} = \frac{1}{2} \left| \det\brak{\begin{bmatrix} \vec{a} & \vec{m} \end{bmatrix}} \right| = \frac{1}{2} \left| \det \myvec{3 & -12/5 \\ 0 & 9/5} \right|
\]
\[
    \text{Area} = \frac{1}{2} \left| \brak{3 \cdot \frac{9}{5}} - \brak{0 \cdot \frac{-12}{5}} \right| = \frac{1}{2} \left| \frac{27}{5} \right| = \frac{27}{10}
\]
The area of the triangle OAM is $\frac{27}{10}$. This corresponds to option (d).
   
\end{frame}

\begin{frame}{}
    \centering
    \includegraphics[width=\columnwidth, height=0.8\textheight, keepaspectratio]{figs/ellipse.png}  
\end{frame}


\end{document}
